% Conclusion. 180-320 words. Only gate-passed claims. No new numbers.
\section{Conclusion}
\label{sec:conclusion}

A comparison between selection strategies for a disagreement-based certificate
is identifiable only when the \splitseed\ and the \trajseed\ are varied
independently under frozen \evalsem. We registered that matrix before running
it, preserved every run with a manifest identifying its method and seeds, and
reported the full set of per-split contrasts rather than the most favorable one.

On MNIST at a reduced probe scale, \methodrandom\ and \methodmax\ are
indistinguishable. The mean paired contrast favours \methodrandom, but it is
smaller than the largest within-split trajectory spread in the matrix, and it
changes sign at one of the five splits. What the matrix does establish is the
size of the nuisance: the trajectory alone moves the certificate by more than
the selection criterion does, and the between-split range grows with how
aggressively a method selects. A study that reports one trajectory per
configuration is measuring that nuisance and attributing it to the method.

The scope is narrow by construction. One dataset, one certificate family, two
selection criteria, and a reduced experimental scale. The comparison at the
upstream scale of 200 epochs and 20{,}000 compression size was not run, and the
scale ablation shows an effect there of the same order as the between-method
contrast, so that comparison is the one most likely to change the answer.
Disagreement-aware and hybrid selection criteria are untested here and need the
same paired design rather than a single additional run. Whether the trajectory
spread narrows at full scale, or is an intrinsic property of the certificate, is
the question this design is now sized to answer.

\section*{Acknowledgments}
This document is a demonstration instance produced by the AutoConference
paper-writing skill. Its numbers are synthetic and describe no real experiment.
